% Clock Pre-Registration v3 (FINAL)
% Merges Claude v2 + i20 Agent findings + referee corrections
\documentclass[aps,prd,twocolumn,superscriptaddress,floatfix]{revtex4-2}
\usepackage{amsmath,amssymb,booktabs,xcolor}
\usepackage[colorlinks=true,linkcolor=blue,citecolor=blue,urlcolor=blue]{hyperref}
\newcommand{\kalpha}{k_\alpha}
\newcommand{\epsH}{\epsilon_H}
\newcommand{\azero}{a_0}
\begin{document}
\title{Pre-Registered Predictions for Multi-Species Atomic Clock Tests\\
of Composition-Dependent Scalar-Field Gravity}
\author{G.~T.~Alcock}
\affiliation{Independent Researcher}
\date{\today}
\begin{abstract}
We pre-register a pattern of predictions for annual modulation of
atomic clock frequency ratios in a scalar-field gravity framework.
The framework predicts a channel-resolved coupling hierarchy with
three layers: a common electromagnetic sector, composition-sensitive
residuals suppressed by $\epsilon_H = 1/20$, and a pure
electromagnetic residual suppressed by $\epsilon_H^2$. The absolute
signal amplitudes depend critically on an unsettled screening
mechanism: with the screening derived in the parent
theory, amplitudes are $\sim\!10^{-20}$--$10^{-22}$ and
undetectable; without screening, they reach $\sim\!10^{-16}$ for
cross-species pairs. A critical exception is the $^{229}$Th nuclear
clock, where strong-sector enhancement ($S^{\alpha_s} \sim 10^4$)
pushes the screened amplitude to $\sim\!10^{-18}$---potentially
detectable. We cannot resolve the screening ambiguity theoretically
at present. However, the \textit{screening-independent} observables
--- inter-pair amplitude ratios, perihelion phase lock, and
same-ion null --- form a four-test discriminating pattern that
no generic scalar-tensor theory reproduces without additional free
parameters. The nearest-term falsification test is the PTB Yb$^+$
E3/E2 bound on $\lambda_\alpha$, which will reach the predicted
value $\lambda_\alpha \approx 2.1 \times 10^{-8}$ within
$\sim$1--2 years. We specify signal templates, blinded analysis
protocol, and five quantitative falsification criteria.
\end{abstract}
\maketitle
%=============================================================================
\section{Introduction}
%=============================================================================
Local Position Invariance (LPI) requires that the gravitational
redshift is universal~\cite{Will2014}. Any species-dependent
deviation constitutes a violation of general relativity. This
paper pre-registers predictions of Density Field
Dynamics~\cite{Alcock2026DFD} (DFD), a scalar-tensor theory with
$n = e^\psi$ on flat spacetime.
\paragraph{Scope.} Clock tests probe DFD's \textit{gauge-coupling
sector} ($\kalpha$, $\lambda_I$, $\epsH$). They are orthogonal
to DFD's core gravitational tests (GW non-lensing, rotation curves,
$D_L^{\rm EM}/D_L^{\rm GW}$ ratio). A positive clock result
establishes composition-dependent gravity and LPI violation---not
DFD as a complete theory. A null constrains only the gauge sector.
%=============================================================================
\section{Theoretical Framework}
%=============================================================================
\subsection{The coupling hierarchy}
For co-located clock ratios $R_{AB} = \nu_A/\nu_B$, the common
gravitational redshift cancels identically. The observable is
\begin{equation}\label{eq:ratio-obs}
\frac{\Delta R_{AB}}{R_{AB}}
= (K_A^{\rm obs} - K_B^{\rm obs})\,\frac{\Delta\Phi}{c^2},
\end{equation}
where the clock coefficient is
\begin{equation}\label{eq:Kobs}
K_A^{\rm obs} = \Sigma\bigl[\lambda_\alpha\,\widetilde{S}_A^\alpha
+ \lambda_N\,C_N^{(A)} + \lambda_e\,C_e^{(A)}\bigr],
\end{equation}
with screening factor $\Sigma$ (discussed below). The hierarchy:
\begin{align}
\kalpha &= \alpha^2/(2\pi) \approx 8.5 \times 10^{-6}
 \quad \text{(common, unobservable in ratios)}, \label{eq:kalpha}\\
\lambda_{N,e} &\approx \epsH\,\kalpha \approx 4.2 \times 10^{-7}
 \quad \text{(composition channel)}, \label{eq:lcomp}\\
\lambda_\alpha &\approx \epsH^2\,\kalpha \approx 2.1 \times 10^{-8}
 \quad \text{(pure-$\alpha$ residual)}, \label{eq:lalpha}
\end{align}
where $\epsH = N_{\rm gen}/k_{\max} = 3/60 = 1/20$.
\textit{Derivation status.}
$\kalpha$ follows from the photon-$\psi$ vertex at one loop
(Schwinger factor $\alpha/(2\pi)$ times atomic Coulomb
$\alpha$)~\cite{Alcock2026DFD}. This requires one axiom beyond
DFD's four (``one gauge vertex''). The integers 3 and 60 are
topological invariants; their ratio as the class-breaking parameter
is well-motivated but not formally proven (grade~B). The power-law
scaling ($\epsH$ vs $\epsH^2$) follows from class-insertion counting.
\subsection{The screening problem}\label{sec:screening}
The screening factor from Ref.~\cite{Alcock2026DFD}, Eq.~(9.14):
\begin{equation}\label{eq:Sigma}
\Sigma(y) = 1/\sqrt{1+y}, \quad y \equiv a/\azero,
\end{equation}
is derived from a postulated response functional, not from the
DFD action. On Earth's surface ($y \approx 8 \times 10^{10}$):
$\Sigma_\oplus \approx 3.5 \times 10^{-6}$.
\textit{The critical consequence:} combining Eqs.~(\ref{eq:lcomp})
and (\ref{eq:Sigma}) with $\Delta\Phi_\odot^{\rm peak}/c^2 \approx
1.65 \times 10^{-10}$:
\begin{equation}
A_{\rm Cs/Sr}^{\rm screened} \sim 4.2\!\times\!10^{-7}\times
3.5\!\times\!10^{-6}\times 1.65\!\times\!10^{-10}
\approx 2\!\times\!10^{-22}.
\end{equation}
This is \textbf{undetectable}. Even at solar-orbit screening
($\Sigma_\odot \approx 1.4 \times 10^{-4}$), the amplitude is
$\sim\!10^{-20}$---still undetectable.
If screening does \textit{not} apply (e.g., if the External Field
Effect prevents local screening of the external solar potential
variation, as in MOND~\cite{Milgrom1983}):
\begin{equation}
A_{\rm Cs/Sr}^{\rm unscreened} \sim 4.2\!\times\!10^{-7}\times
3\times 1.65\!\times\!10^{-10} \approx 2\!\times\!10^{-16}.
\end{equation}
This would be detectable.
\textbf{We cannot resolve this ambiguity.} The response functional
is model-grade. The EFE analogy is physically motivated but
unproven for DFD's specific screening. A first-principles
derivation from the DFD action is needed.
%=============================================================================
\section{What We Pre-Register}
%=============================================================================
Given the screening ambiguity, we pre-register
\textit{screening-independent observables}: quantities in which
$\Sigma$ cancels.
\subsection{Primary: inter-pair amplitude ratios}
For any two cross-species pairs, the screening cancels:
\begin{equation}
\frac{A_{\rm Yb^+(E3)/Sr}}{A_{\rm Cs/Sr}}
\approx \frac{|\Delta S_{\rm Yb^+(E3)/Sr}|}
{|\Delta S_{\rm Cs/Sr}|}
= \frac{6.01}{2.77} = 2.17,\ \text{opposite sign.}
\end{equation}
This ratio is fixed by atomic structure calculations
($S_A^\alpha$ from Refs.~\cite{Flambaum2009,Dzuba2018}), the
$\epsH$ hierarchy, and the composition coefficients. It contains
no free parameters and is independent of $\Sigma$, $\kalpha$,
and the overall coupling scale.
A Damour-Polyakov model~\cite{Damour1994} has independent dilaton
charges that can match any single ratio. DFD fixes all ratios
through $\epsH$. Two simultaneous pairs constitute a
model-discriminating test.
\subsection{Secondary: phase lock and channel ordering}
\begin{enumerate}
\item \textbf{Perihelion phase.} All signals peak within $\pm 45$
  days of January~3 (maximum solar potential at Earth).
\item \textbf{Same-ion null.} Yb$^+$ E3/E2 is predicted $\sim\!20\times$
  smaller than cross-species pairs (the $\epsH^2/\epsH$ gap).
\item \textbf{Channel ordering.} Cross-species $>$ same-ion. Microwave-optical
  (Cs/Sr) comparable to ion-neutral (Yb$^+$/Sr).
\end{enumerate}
\subsection{The nuclear clock channel: $^{229}$Th}\label{sec:nuclear}
The full observable coupling [Eq.~(9.5) of
Ref.~\cite{Alcock2026DFD}] includes a strong-sector term
$\lambda_s\,S_A^{\alpha_s}$ omitted from Eq.~(\ref{eq:Kobs})
because it vanishes for electronic transitions. For the
$^{229}$Th nuclear isomer ($E \approx 8.34$~eV~\cite{Tiedau2024}),
the strong-sector sensitivity is
$S^{\alpha_s} \sim 10^4$~\cite{Flambaum2006nuclear}, four orders
of magnitude above any electronic transition.
Within the class-breaking hierarchy,
$\lambda_s \approx \epsH\,\kalpha \approx 4.2 \times 10^{-7}$.
The nuclear-channel annual amplitude is:
\begin{equation}
A_{\rm Th/Sr} \sim \lambda_s\,S^{\alpha_s}\,\Sigma\,
\frac{\Delta\Phi}{c^2}
\sim 4.2\!\times\!10^{-3}\times\Sigma\times
1.65\!\times\!10^{-10}.
\end{equation}
Under Earth-surface screening ($\Sigma \approx 3.5 \times 10^{-6}$):
$A^{\rm screened} \sim 2 \times 10^{-18}$, corresponding to
$\delta\nu_{\rm Th} \approx 5$~Hz---slightly below the Ooi
et al.~\cite{Ooi2026} reproducibility floor of 26~Hz. Under
solar-orbit screening: $\sim\!10^{-16}$ ($\delta\nu_{\rm Th}
\approx 140$~Hz), well within the surviving window
($26$~Hz--$\mathcal{O}(1~\rm kHz)$). Unscreened: $\sim\!10^{-13}$
(already excluded). The Th-229 channel thus helps resolve the
screening ambiguity: detection in the 26--1000~Hz window favors
solar-orbit screening; a null below 26~Hz is consistent with
Earth-surface screening.
\textbf{This is potentially the only DFD clock channel detectable
even with full Earth-surface screening.}
The nuclear-to-optical amplitude ratio is screening-independent:
\begin{equation}\label{eq:nuclear-ratio}
\frac{A_{\rm Th/Sr}}{A_{\rm Cs/Sr}} \approx
\frac{\lambda_s\,S^{\alpha_s}}{\lambda_N\,\Delta C_N}
\sim \frac{10^4}{3} \approx 3000.
\end{equation}
This is the single most discriminating pre-registered prediction:
if annual modulation is detected in any pair, the nuclear-to-optical
ratio sharply tests the strong-sector coupling hierarchy.
\subsection{Nearest-term test: PTB E3/E2}
DFD predicts $\lambda_\alpha \approx 2.1 \times 10^{-8}$
(Eq.~\ref{eq:lalpha}). The PTB bound is $|\lambda_\alpha|
\lesssim 3.2 \times 10^{-8}$ ($2\sigma$)~\cite{Filzinger2023}.
The prediction sits at $66\%$ of the bound. A factor-of-2
improvement (expected $\sim$1--2 years) is decisive: either
detection or exclusion of the $\epsH^2$ prediction. This test
is \textit{screening-independent} because E3/E2 probes a static
potential, not an annual modulation.
%=============================================================================
\section{Signal Template and Analysis Protocol}
%=============================================================================
The annual signal for pair $A/B$ is
\begin{equation}
\frac{\Delta R_{AB}(t)}{R_{AB}}
= A_{AB}\cos[M(t)] + \beta_0 + \beta_1 t + \epsilon(t),
\end{equation}
where $M(t)$ is Earth's mean anomaly ($M = 0$ at perihelion).
\paragraph{Blinding.} A random offset is added to the
perihelion-phase amplitude before analysis. The pipeline
(regression model, uncertainty estimation, phase tests) is
frozen before unblinding.
\paragraph{Phase-specificity test.} The regression is repeated
with four alternative phase hypotheses (aphelion, spring
equinox, summer solstice, fall equinox). A DFD signal requires
perihelion amplitude $> 3\sigma$ with all alternatives $< 2\sigma$.
\paragraph{Uncertainty estimation.} Least-squares with analytical
errors, validated by leave-one-day-out jackknife, wild bootstrap,
and sign-permutation resampling.
%=============================================================================
\section{Falsification Criteria}
%=============================================================================
\begin{enumerate}
\item \textbf{Same-ion signal:} Detection in Yb$^+$ E3/E2 at
  amplitude within $3\times$ of cross-species pairs falsifies
  the $\epsH$ hierarchy.
\item \textbf{Wrong ratio:} Yb$^+$(E3)/Sr to Cs/Sr deviates from
  $2.17 \pm 0.8$ by $> 3\sigma$.
\item \textbf{Wrong phase:} Any cross-species signal peaks
  $> 45$ days from perihelion at $> 3\sigma$.
\item \textbf{Wrong sign:} Yb$^+$(E3)/Sr and Cs/Sr have the
  same sign.
\item \textbf{Universal null:} All cross-species ratios null at
  $< 10^{-18}$ with $> 5\sigma$ across $\geq 3$ pairs
  (excludes unscreened scenario).
\item \textbf{PTB exclusion:} $|\lambda_\alpha| < 1.0 \times 10^{-8}$
  at $2\sigma$ excludes the $\epsH^2$ prediction (partial
  falsification of the class-insertion hierarchy).
\item \textbf{Nuclear ratio:} If Th/Sr and Cs/Sr are both detected,
  their ratio deviating from $\sim\!3{,}000$--$10{,}000$
  (depending on $|\Delta C_N|$) by $> 3\sigma$
  excludes the strong-sector hierarchy
  (Eq.~\ref{eq:nuclear-ratio}). A Th/Sr amplitude outside the
  surviving window ($26$~Hz--$\mathcal{O}(1~\rm kHz)$) excludes
  the screened strong-sector prediction.
\end{enumerate}
%=============================================================================
\section{The ROCIT Anomaly}
%=============================================================================
The ROCIT reanalysis~\cite{Alcock2025_ROCIT} reports a $13.5\sigma$
perihelion-locked signal in Yb$^+$(E3)/Sr at $A = (-1.045 \pm
0.078) \times 10^{-17}$. This amplitude is 5--7 orders above the
screened predictions and $\sim\!10\times$ below the unscreened
predictions. Three interpretations exist: (i)~systematic artifact,
(ii)~DFD with intermediate screening ($\mathcal{S} \approx 14$),
(iii)~non-DFD new physics. This protocol is designed to confirm
or refute it with independent data.
If ROCIT is confirmed, the ratio test becomes immediately
decisive: does the Cs/Sr amplitude equal $A_{\rm ROCIT}/2.17$
with opposite sign? If yes, the $\epsH$ hierarchy is supported.
If not, the channel structure is excluded.
%=============================================================================
\section{Discussion}
%=============================================================================
\paragraph{Honest assessment.} DFD's clock sector is theoretically
structured but experimentally ambiguous. The coupling hierarchy
($\kalpha$, $\lambda_I$, $\epsH$) produces specific testable
ratios, but the screening mechanism needed for consistency with
existing null results (BACON~\cite{BACON2021}, PTB) also
suppresses predicted signals below detectability. The theory
makes sharp predictions for the \textit{pattern} while leaving
the \textit{overall scale} uncertain.
\paragraph{Three paths forward.}
(i)~\textit{Derive screening from first principles} (the DFD path
integral, not a postulated free energy). This is a well-defined
calculation.
(ii)~\textit{Run PTB E3/E2 to completion} ($\sim$1--2 years).
This tests $\lambda_\alpha$ directly, independent of screening.
(iii)~\textit{Space-based clocks} in low-acceleration environments
($a \sim 10^{-6}$~m/s$^2$, $\Sigma \sim 10^{-2}$) would increase
sensitivity by four orders of magnitude.
(iv)~\textit{Nuclear clock Th/Sr comparison} over $\geq 1$ year:
the $10^4$ strong-sector enhancement may overcome screening,
placing the signal at $\sim\!10^{-18}$---within reach of
next-generation nuclear clocks now becoming operational.
\paragraph{Comparison with Damour-Polyakov.}
The distinctive DFD features are: (i)~$\kalpha = \alpha^2/(2\pi)$
(a specific numerical value, not free), (ii)~$\lambda_\alpha :
\lambda_N = 1 : 20$ (the $\epsH$ ratio), (iii)~perihelion phase
(solar potential driver), (iv)~$\Sigma = 1/\sqrt{1+y}$ (distinct
from chameleon or symmetron). Simultaneous detection matching all
four would be strong evidence for DFD specifically.
%=============================================================================
\section{Conclusion}
%=============================================================================
We pre-register a \textit{pattern}, not an amplitude. The absolute
signal strength depends on a screening mechanism not yet derived
from first principles; under current screening, signals are
undetectable. But the screening-independent pattern---fixed
inter-pair ratios, perihelion phase, same-ion null, and the 1:20
channel hierarchy---is unique to DFD and falsifiable by multiple
quantitative criteria. The nearest-term test is PTB Yb$^+$ E3/E2,
expected decisive within $\sim$1--2 years. The $^{229}$Th nuclear
clock may be the only channel where screening does not prevent
detection, with a predicted nuclear-to-optical ratio of $\sim\!3000$
that is screening-independent. If confirmed, the ratio test across
multiple pairs becomes the definitive discriminant.
\begin{acknowledgments}
The author thanks the anonymous referees whose audits identified
the screening inconsistency and derivation-chain gaps in v1,
leading to the honest two-scenario presentation adopted here.
\end{acknowledgments}
\begin{thebibliography}{99}
\bibitem{Will2014} C.~M.~Will,
\textit{Living Rev.\ Relativ.}\ \textbf{17}, 4 (2014).
\bibitem{Alcock2026DFD} G.~T.~Alcock, Zenodo (2026),
\url{https://zenodo.org/records/19029160}.
\bibitem{Milgrom1983} M.~Milgrom,
\textit{Astrophys.\ J.}\ \textbf{270}, 365 (1983).
\bibitem{Damour1994} T.~Damour and A.~M.~Polyakov,
\textit{Nucl.\ Phys.\ B}\ \textbf{423}, 532 (1994).
\bibitem{Flambaum2009} V.~V.~Flambaum and V.~A.~Dzuba,
\textit{Can.\ J.\ Phys.}\ \textbf{87}, 25 (2009).
\bibitem{Dzuba2018} V.~A.~Dzuba, V.~V.~Flambaum, and S.~Schiller,
\textit{Phys.\ Rev.\ A}\ \textbf{98}, 022501 (2018).
\bibitem{Filzinger2023} M.~Filzinger \textit{et al.},
\textit{Phys.\ Rev.\ Lett.}\ \textbf{130}, 253001 (2023).
\bibitem{Alcock2025_ROCIT} G.~T.~Alcock, (2025),
archived at Ref.~\cite{Alcock2026DFD}.
\bibitem{BACON2021} BACON Collaboration,
\textit{Nature}\ \textbf{591}, 564 (2021).
\bibitem{Tiedau2024} J.~Tiedau \textit{et al.},
\textit{Phys.\ Rev.\ Lett.}\ \textbf{132}, 182501 (2024).
\bibitem{Flambaum2006nuclear} V.~V.~Flambaum,
\textit{Phys.\ Rev.\ Lett.}\ \textbf{97}, 092502 (2006).
\bibitem{Ooi2026} T.~Ooi \textit{et al.},
``Reproducibility of $^{229}$Th nuclear clock transitions,''
(2026).
\end{thebibliography}
\end{document}
